\label{sec:dataset}
\paragraph{Challenges and overview.}
Prior studies on building LLMs to synthesize scientific literature employ either small-scale, single-domain human evaluation~\citep{agarwal2024litllm,zheng2024openresearcher} or over-simplified multiple-choice QA setups~\citep{skarlinski2024language}. 
Building high-quality benchmarks for literature review has two major challenges. First, creating such datasets is resource-intensive, as it requires Ph.D.-level domain expertise and research experience, particularly when annotating realistic questions and high-quality answers. Second, even when high-quality data is available, reliably evaluating long-form natural language responses presents a significant challenge, especially in expert domains~\citep{xu2023critical,xu-etal-2024-kiwi}. This contrasts with benchmarks for other scientific processes, such as automated experimental code generation, for which clearer evaluation criteria, such as Pass@1, are more readily available~\citep{si2024can}.

To address these gaps, we introduce \data, a  benchmark that supports diverse formats of scientific literature synthesis tasks, including closed-form classification, multiple-choice, and long-form generation, as shown in Table~\ref{table:overview_dataset}.
We adopt three existing single-paper datasets, and then construct a suite of high-quality, expert annotated datasets for computer science, biomedicine,  physics, and neuroscience (Section~\ref{subsec:data_curation}). We also build a reliable automatic evaluation pipeline (Section~\ref{subsec:automatic_eval}). 
\Tref{table:overview_dataset} provides a list of tasks in \data, and Figure~\ref{fig:overview_multi_paper_qa} shows an example and an overview of the evaluation pipeline. 


\begin{table}[t]
\small
    \centering
    \resizebox{\textwidth}{!}{
    \begin{tabular}{c  ccccc}
    \toprule
        \textbf{Dataset} & \textbf{Task Format} & \bf Discipline & \bf Size & \bf Evaluation & \bf Multi-paper  \\
    \midrule
    SciFact & Claim $\rightarrow$ Label & Biomedicine & 208 & \acc, \cit  \\
    \textcolor{gray}{\footnotesize{(\citealt{wadden-etal-2020-fact})}} & (True or False) &  \\\midrule
    PubMed QA & Question $\rightarrow$ Answer  & Biomedicine & 843 & \acc, \cit  \\
   \textcolor{gray}{\footnotesize{(\citealt{jin-etal-2019-pubmedqa})}} & (Yes, No) & \\ \midrule
    % BioASQ & \\
    QASA & Question $\rightarrow$ Answer  & Computer Science & 1,375 & \acc, \cit & \\
    \textcolor{gray}{\footnotesize{(\citealt{lee2023qasa})}} & (Long-form) & \\ \midrule\midrule
    \textsc{ScholarQA-CS}& Question $\rightarrow$ Answer$^\dagger$ & Computer Science & 100 & \acc,\cit  & \multirow{2}{*}{\checkmark}\\
    & (Long-form) &  & &  \\ \midrule
    \textsc{ScholarQA-Bio} & Question $\rightarrow$ Answer$^*$ & Biomedicine & 1,451 & \cit  & \multirow{2}{*}{\checkmark}\\
    & (Long-form) &  & & \\ \midrule
    \textsc{ScholarQA-Neuro} & Question $\rightarrow$ Answer$^*$ & Neuroscience & 1,308 & \cit  & \multirow{2}{*}{\checkmark}\\
    & (Long-form) &  & & \\  \midrule
    \textsc{ScholarQA-Multi} & Question $\rightarrow$ Answer & Computer Science, Physics, & 108 & \cit,  & \multirow{2}{*}{\checkmark}\\
    & (Long-form) &  Biomedicine &   & \llm, \expert \\
    \bottomrule
    \end{tabular}
    }
        \caption{{\bf Overview of \data}. The top three rows show single-paper datasets adopted from prior datasets. 
        The bottom four rows are new datasets, which we constructed by recruiting Ph.D.-level experts. 
        Answer$^*$ indicates that the dataset comes with questions only, and answer$^\dagger$ indicates that answer will be evaluated based on human-annotated rubrics. 
        The evaluation columns corresponds to the multi-faceted evaluations in Section~\ref{subsec:automatic_eval}. The ``Multi-paper'' column indicates whether the task requires multiple papers to answer. 
        \llm~and \expert~indicate the fine-grained evaluations by evaluator LLMs (i.e., Prometheus; \citealt{kim2024prometheus}) and expert humans, respectively. 
    }
    \label{table:overview_dataset}
\end{table}


\subsection{Data Curation}
\label{subsec:data_curation}
\data is designed to evaluate model capabilities in automating scientific literature review. The curation process is guided by three key factors: 
{\bf Diversity of tasks}: \data includes tasks with a range of input-output formats; 
{\bf Diversity of disciplines}: Unlike previous analyses that often focus on a single discipline such as computer science, 
\data spans four scientific disciplines; 
{\bf Inclusion of multi-paper tasks}: Unlike prior work that focuses on understanding single, pre-selected papers, all tasks require retrieving from the entire open-access collection of full text of papers (\Sref{subsec:single_paper_task}) 
 and four datasets specifically require reasoning over multiple retrieved papers (\Sref{subsec:multi_papers}). 


\subsubsection{Single-Paper Tasks}
\label{subsec:single_paper_task}
For single-paper tasks, we curate and adapt existing widely-used single-paper datasets. Figure\ref{fig:single_paper_examples} shows examples of single-paper tasks; more details are in Appendix~\ref{app_sec:annotation_details}. 

\paragraph{SciFact.}~SciFact \citep{wadden-etal-2020-fact} is a dataset of 1.4K expert-written scientific claims in the biomedical domain, paired with gold evidence from existing PubMed paper abstracts annotated with labels and rationales. 
We include validation set queries labeled as either \texttt{supports} (true) or \texttt{contradicts} (false), discarding the original gold evidence, and reformulate the task as binary open-retrieval, wherein a system needs to identify relevant papers from a large collection of papers.  

\paragraph{PubMedQA.}
{PubMedQA}~\citep{jin-etal-2019-pubmedqa} has expert-annotated (yes/no/maybe) QA data on PubMed paper abstracts. 
Similarly to {SciFact}, we only keep instances with yes or no labels, and discard the original abstract passage to formulate the task as an open-retrieval setup. 

\paragraph{QASA.}
{QASA}~\citep{lee2023qasa} is a single paper QA dataset that consists of question answering pairs, requiring reasoning over scientific articles in AI and ML. 
We evaluate the model's ability to sufficiently answer a detailed question about the target paper. While the original dataset provides three subtasks (answer selection, rationale generation and answer compositions) as well as end-to-end QA, we evaluate the models' performance based on an end-to-end QA setup. 




\begin{figure*}[t!]
    \centering
    \includegraphics[width=\textwidth]{figs/scilit_example_figure_v2.pdf}
    \caption{{\bf An \textsc{ScholarQA-CS} example and evaluation overview.} \textsc{ScholarQA-CS} consists of 100 questions and an average of 4.4 expert-written rubrics to be satisfied. Our \data evaluation pipeline evaluates aspects like correctness and citation accuracy.  
    }
    \label{fig:overview_multi_paper_qa}
\end{figure*}

\subsubsection{Multi-paper Tasks}
\label{subsec:multi_papers}
Single-paper, closed-set tasks may provide reliable evaluations. However, they may not be reflective of realistic scenarios, in which complex, open-ended questions are asked independently from existing papers, and require multi-paper retrieval and reasoning. 
Few datasets~\citep{xu-etal-2024-kiwi,malaviya2023expertqa} explore multi-paper setups with realistic queries, and most lack a reliable evaluation pipeline or human-written references. We address this gap by curating three new long-form QA datasets, annotated by experts, for these challenging settings (details in Appendix~\ref{app_sec:annotation_details}). 
Furthermore, our multi-paper tasks include four scientific disciplines. 



\paragraph{\textsc{ScholarQA-CS}.} 
We collected 100 questions along with detailed answer rubrics for each question across various computer science disciplines by recruiting expert annotators holding Ph.D.s in the field (professors, postdoctoral researchers, and research scientists). 
Annotators were tasked with writing literature review questions that were expected to require multiple research papers to answer. 
The question topics span areas such as networks, algorithms, the Internet of Things, artificial intelligence, and human-computer interaction. 
Then, for each question, two other annotators searched the Web to produce a rubric listing the key ingredients for a correct answer, categorized by importance (``must-have'' and ``nice-to-have''), along with supporting quotes from sources for each ingredient. {The annotators were instructed not to use any LLM services for this initial part of the task. After the initial web search, the annotators were shown corresponding responses from four LLM services {(Claude 3.5 Sonnet, GPT-4o, Perplexity Pro and an unpublished RAG prototype based on Claude 3.5)} in a randomized order in case they wanted to revise their rubrics.} On average, each question is annotated with 4.4 key ingredients, each supported by 4.4 quotes. 

{To measure agreement, we had both annotators produce rubrics for a subset of 10 randomly sampled questions. We then compute the scores for responses from the four LLM services the annotators were exposed to using our automated approach, once for each set of annotator rubrics. Finally, we compute Pearson's correlation coefficient among the scores for each question and take the mean. Since the rubric annotation task is subjective, we compute this agreement both with and without the general criterion as part of the scores, which comes out to be 79.3 and 59.5, respectively.}
Figure~\ref{fig:overview_multi_paper_qa} shows one example, and more examples and details are available in Appendix \ref{app_sec:examples_rubric_qa}. 

\paragraph{\textsc{ScholarQA-Bio}, \textsc{ScholarQA-Neuro}.}
We further collected 2,759 expert-written literature review questions in biomedicine and neuroscience, recruiting six experts who have a Ph.D. in relevant areas and are currently research scientists and engineers.  
{The annotators were asked to choose papers from their area of expertise,} and generate complex scientific questions that biomedical scientists might reasonably ask about the scientific literature {based upon their parsing of those papers}. We collected questions from different areas, such as bioimaging, genetics, microbiology, and neuromodulation for each.  
Due to the cost of annotation, we focused solely on curating the questions. 
Full instructions and examples are available in Appendix~\ref{tab:annotator instructions_biomed} and \ref{app_sec:examples_biomed_queries}. 

\paragraph{\textsc{ScholarQA-Multi}.}
Lastly, we collected 108 literature review questions and expert-written answers with citations in four domains: computer science (AI/ML, HCI), Biomedicine (Bioimaging,  Genetics), Physics (Astrophysics, Photonics, Bio Physics). 
All annotations are conducted by Ph.D. students or post-Ph.D. scientists, who have more than three years of research experience in the corresponding areas and have multiple first-author publications. 
We asked them to come up with questions that are related to most recent literature, and to compose answers to the questions using relevant papers that they found via search. Our annotators were instructed not to use any LLM-based systems such as ChatGPT, and told to only use general search (e.g., Google Search) or paper search systems (e.g., Semantic Scholar). 
Table~\ref{table:result_human_data} show statistics of the collected questions and answers and the distribution of subjects are in Figure~\ref{fig:data_fistributions}, along with the average annotation time per subject. We show several examples in Appendix \ref{app_sec:examples_annotated}. 
On average, each annotator spent 56 minutes per instance.

\subsection{Metrics and Evaluation Protocols}
\label{subsec:automatic_eval}
We developed a multifaceted automatic evaluation pipeline to facilitate reproducible and efficient evaluations, complementing expert assessments. An overview of our evaluations is in Figure~\ref{fig:overview_multi_paper_qa}. 


\paragraph{Correctness (\acc).}
Correctness evaluates the degree of overlap or matching of a model-generated answer and human-annotated reference. This metric is only applied to tasks with human-annotated reference answers. 
For short-form generation tasks given a fixed set of answer classes, namely SciFact and PubMedQA, we use accuracy as the correctness metric. 
For QASA, we use ROUGE-L as an evaluation metric, following \cite{lee2023qasa}. 
For \textsc{ScholarQA-CS}, we develop a new long-form evaluation pipeline, which employs expert-annotated rubrics. Each rubric has two criteria: general (accounting for 40\% of the score) and annotation-driven (60\%). 
General criteria cover the evaluation of length, expertise, citations, and excerpts, while annotation-driven criteria involve assessing the presence of specific key ingredients identified by annotators. GPT4o-turbo assigns scores for each criterion, and a weighted sum is computed to obtain a final score. 
More details are in appendix \ref{app_sec:details_scholar_cs}. 

\paragraph{Citation accuracy (\cit).}
Evaluating long-form responses to literature review questions requires citation accuracy: LMs should correctly attribute relevant evidence for all citation-worthy statements. In \data, all systems generate outputs with reference numbers (e.g., \texttt{[1], [2]}) linked to passages provided during inference. Following prior work~\citep{gao2023enabling,liu2023evaluating}, we measure citation precision and recall. Specifically, we check if each citation-worthy statement has appropriate citations and if the citations support the statement ({\bf Citation Recall}, \cit-r). 
For each citation, we then verify its relevance and necessity---specifically, whether the citation supports the statement and if its removal impacts the integrity of remaining citations ({\bf Citation Precision}, \cit-p). 
Finally, we compute {\bf Citation F1} (\cit-F1) as well, and use it as a primarily metric for citation accuracy. 
Citation accuracy does not require gold reference answers or rubrics, so we apply this evaluation across all tasks. 

\paragraph{Content quality and organization (\llm, \expert).}
We further define key aspects to evaluate the generated answers beyond \acc~ or \cit~ alone. Specifically, we assess {\bf Relevance} (\rel) to the question, {\bf Coverage} (\cov) of topics (e.g., diversity of discussed papers) and depth (e.g., sufficiency of details), and {\bf Organization and Writing Flow} (\org). These aspects are challenging to capture with standard metrics.  
 Since LMs can effectively follow detailed evaluation rubrics~\citep{zheng2023judging,kim2024prometheus}, we use Prometheus v2~\citep{kim2024prometheus} to assign five-scale scores based on defined rubrics and use the same schema for human evaluations. For human evaluation, we further evaluate {\bf Overall Usefulness}  (\useh).  
Full instructions for this evaluation are provided in the Appendix \ref{app_section:evaluations}. 
As prior studies show that \llm~is less reliable when gold reference answers are not available~\citep{kim2024biggen}, this evaluation is only applied to a task with human-annotated reference answer, namely \textsc{ScholarQA-Multi}. 
We analyzed the agreement between human and model assessments on fine-grained aspects (Appendix \ref{app_sec:human_evaluation}) and found that the model's evaluations often align with human rankings, showing higher correlation especially in organization and coverage. 
